\section{Simultaneous path pruning and affine witness patterns}\label{sec:patterns}

We now organize the local group arguments into a finite pruning construction on the test lattices of \Cref{sec:affine}. The construction is simultaneous over all ancestors of a path. Its output will be used with the classical and alternating exclusions in \Cref{sec:classical,sec:alternating}, \Cref{w:bounded-exclusion}, and \Cref{w:path-bound}.

\subsection{Path labels and compatible simple groups}

We work in the ordinary interval supplied by
\Cref{f:thm-almost-simple-reduction}, with fixed base group $H$.
Write $K_X$ for the subgroup labeled by a reverse flat $X$. Thus
$X\subseteq Y$ means $K_Y\le K_X$. A downward path chooses, at its current
flat $W$, a private trunk chart and then a flat $U\supsetneq W$ of rank
increment between $1$ and $r$ in that chart. The local group arguments
will be applied both in the current group's almost simple quotient and
in the quotients belonging to all earlier path nodes.

At a path node $W$ we maintain an actual nonabelian simple group
$B_W\trianglelefteq K_W$ such that
\[
 HB_W=K_W,\qquad C_W:=C_{K_W}(B_W)\le H,
\]
where $C_W$ is soluble and $K_W/C_W$ is almost simple with socle the
faithful image of $B_W$. The starting interval has this property. The
next lemma explains why the local conclusions in several such quotients
give one subgroup in the original interval.

\begin{lemma}[Compatibility of ancestor socles]\label{c:ancestor-socles}
Suppose $H\le K\le K_W$, and a local conclusion in $K_W/C_W$ supplies
a nonabelian simple normal subgroup $\overline B\trianglelefteq K/C_W$
contained in the image of $B_W$. Suppose also that
$(H/C_W)\overline B=K/C_W$, and that its centralizer in $K/C_W$ is
soluble and contained in $H/C_W$. Then $\overline B$ has a unique lift
$B\le B_W$, with
\[
 B\trianglelefteq K,\qquad HB=K,\qquad C_K(B)\le H,
\]
and $C_K(B)$ is soluble. The quotient by this centralizer is almost simple.
If the earlier actual socles contain $B_W$, the image of $B$ in every
earlier ancestor quotient agrees with any nonabelian simple normal
subgroup supplied there with soluble centralizer. All the simultaneous
structural conclusions therefore concern this same $B$.

Whenever two comparable candidate labels have these simple and
centralizer properties, their actual simple groups are nested, strictly
if the labels are distinct.
\end{lemma}

\begin{proof}
The quotient map is faithful on $B_W$, since $B_W\cap C_W=1$.
Its inverse on the image of $B_W$ defines $B$, which lies in $K$ and is
normal there. The supplement identity lifts because $C_W\le H$.
A lift of an element centralizing $\overline B$ centralizes $B$: its
commutators with $B$ lie both in $B$ and in $C_W$, so are trivial.
Hence the new centralizer is the inverse image of the quotient
centralizer and is soluble and contained in $H$. Quotienting by it gives
a faithful automorphism group of $B$ containing its inner automorphisms.

For an earlier node $W'$, the inclusion $B\le B_W\le B_{W'}$ makes its
image faithful in $K/C_{W'}$. The same commutator argument identifies
its centralizer there with $C_K(B)/C_{W'}$, which is soluble. Distinct
nonabelian simple normal subgroups would commute, contradicting that
solubility. Thus the simple supplied relative to $W'$ must be this image.
All these quotient maps preserve the exact interval, since their kernels
lie in $H$.

For comparable candidates with $K_U\le K_L$, the simple $B_U$ acts on
$B_L$. Schreier's Theorem makes its outer action trivial, so
$B_U\le B_L C_{K_L}(B_L)$. This is a direct product because $B_L$ is
centerless. Projection of $B_U$ onto the soluble centralizer is trivial,
giving $B_U\le B_L$. If the labels differ, their simples differ because
$HB_U=K_U$ and $HB_L=K_L$.
\end{proof}

\subsection{One pruning round and its failure witnesses}

A structural exclusion forbids containment in specified proper full
stabilizers in an ancestor's almost simple quotient. Such a full
stabilizer is pulled back to an actual subgroup of the original interval.
Avoiding it is preserved when a tested subgroup is enlarged.

At the current path flat $W$, fix one category of these stabilizers,
allowing finitely many ancestor and case choices. A round tests flats
$U\supsetneq W$ of a prescribed rank increment $s$ in each available
private chart. A successful flat has a label $K_U$ contained in none of
the stabilizers under test. After the first round, tests take place
inside a previously selected successful flat $V$, with
$s<\operatorname{rk}V-\operatorname{rk}W$. The subgroup $E=K_V$ then
satisfies $H<E<K_U$ and retains all preceding exclusions. We call it the
\emph{anchor} for that chart.

\begin{lemma}[Prepared full-stabilizer witnesses]\label{c:prepared-witness}
Suppose $K_W$ already avoids the final exclusion list relative to every
strict ancestor. A failed test at $U$ supplies an actual label $K_X$ with
\[
 W'\subseteq X\subseteq U,\qquad X\not\subseteq W,
\]
where $W'$ is the ancestor at which the failure occurs. If all trial
ranks lie below $m$, this witness has a prepared reverse lower interval.
In a round with anchor $E$, it also satisfies $H<E<K_X$.
\end{lemma}

\begin{proof}
The proper full stabilizer contains $K_U$, so its label lies between
$K_U$ and $K_{W'}$. It is therefore a subflat $X$ of $U$ containing $W'$.
If $X\subseteq W$, its subgroup contains $K_W$, contrary to the earlier
exclusions; for $W'=W$, properness gives the same conclusion.
Such an $X$ is nonempty and lies in the trunk chart containing $U$, so
\Cref{f:prop-prepared-links} applies. Finally $E<K_U\le K_X$.
\end{proof}

\subsection{Patterns forced by a failed round}

\begin{lemma}[Affine patterns in failed charts]\label{c:failed-pattern}
Suppose every increment-$s$ test in a sufficiently large available flat
fails, all trial ranks lie below $m$, and a witness
$X(U)\subseteq U$, with $X(U)\not\subseteq W$, is chosen for each test.
Color it by its ancestor, category and case. Given any finite number of
parameter directions, sufficient available dimension produces a family
$X_t$, indexed by an affine space over $\mathbb F_q$, with these properties:
\begin{enumerate}[label=\textup{(\roman*)}]
\item All witnesses have the same color and rank. For distinct $t,t'$,
their span has rank one larger, contains every $X_u$ on the parameter
line through $t,t'$, and gives the subgroup intersection
$K_{X_t}\cap K_{X_{t'}}$. Each witness label covers this intersection.
\item An affinely independent parameter list increases the span rank by
one at each new point, giving strictly descending subgroup intersections
for as long as the spans remain within the truncation.
\item In an anchored round every label contains the same $E>H$.
Families in different private charts have pairwise subgroup
intersections exactly $H$.
\end{enumerate}
The first round can supply more labels than the full interval height
plus any fixed constant, without requiring their total span to be proper.
\end{lemma}

\begin{proof}
Choose a fixed model for the horizontal increment and many unused
vertical coordinates, with coordinate space $V_{\rm vert}$.
For nonempty $W$, take an origin in $W$ and graph
linear maps from its $s$ new horizontal directions to the vertical
directions; these maps vanish on $W$. For empty $W$, use $s-1$ horizontal
affine dimensions and graph affine maps instead. Every graph is an
increment-$s$ test flat.

Color each map by its witness's failure color and the pullback of $X(U)$
into the fixed horizontal model. There are finitely many colors for the
fixed horizontal data. View the maps as words whose letters are
individual-coordinate linear or affine functionals. Multidimensional
Hales--Jewett \citep{HalesJewett1963} gives a monochromatic family with
fixed coordinates and any prescribed number $k$ of disjoint nonempty
variable batches. Thus the pulled-back witness is one fixed flat
$X_{\rm hor}\not\subseteq W$.

Choose a horizontal functional $\lambda$ nonzero at some point of $X_{\rm hor}$;
in the nonempty case it vanishes on $W$. Restrict the functional in each
variable batch to a scalar multiple of $\lambda$. Let
$\iota:\mathbb F_q^k\to V_{\rm vert}$ send each parameter coordinate
to the sum of the vertical basis vectors in its batch. The disjoint
nonempty batches make $\iota$ injective. The witness family is
\begin{equation}\label{c:graph-witnesses}
 X_t=\{(x,f_0(x)+\lambda(x)\iota(t)):x\in X_{\rm hor}\},
 \qquad t\in\mathbb F_q^k,
\end{equation}
where $f_0$ vanishes on the variable batches and records the fixed graph
coordinates elsewhere. At a point where $\lambda\ne0$, distinct parameters
have different vertical coordinates, so the witness flats are distinct.
The span of two graph flats adds exactly the vertical direction
$\iota(t'-t)$, and therefore contains the witnesses on the parameter
line. More generally,
\begin{equation}\label{c:graph-span-rank}
 \operatorname{rk}\langle X_{t_0},\ldots,X_{t_h}\rangle
 =\operatorname{rk}X_{\rm hor}+
   \dim\operatorname{aff}\{t_0,\ldots,t_h\}.
\end{equation}
The chart join rule converts these spans, when proper, into the asserted
subgroup intersections and covers. Pair spans are proper because the
trial ranks are below $m<n-2$. The anchor inclusion follows from the
round setup. Across private charts the witness flats lie outside $W$
and have no common proper primal upper bound, so their reverse meet is
the base $H$.

Taking sufficiently many parameter directions supplies the first-round
family of more labels than the interval height. Only pair spans and the
bounded independent lists used later must fit the truncation; the span
of the whole family need not be proper. In subsequent rounds the number
of parameter directions is bounded independently of $n$.
\end{proof}

\begin{figure}[tbp]
\centering
\begin{tikzpicture}[x=1cm,y=1cm,
  witness/.style={font=\small,inner sep=2pt},
  inclusion/.style={densely dashed},
  panel/.style={draw=black!30,rounded corners=3pt,inner sep=7pt}]
\node[witness] (h) at (0,0) {\(H\)};
\node[witness] (t1) at (-2.4,1.05) {\(T_1\)};
\node[witness] (t2) at (2.4,1.05) {\(T_2\)};
\node[witness] (k11) at (-3.7,2.2) {\(K_{X_1}\)};
\node[witness] (k12) at (-2.4,2.2) {\(K_{X_2}\)};
\node[witness] (k1q) at (-1.1,2.2) {\(K_{X_q}\)};
\node[witness] (k21) at (1.1,2.2) {\(K_{Y_1}\)};
\node[witness] (k22) at (2.4,2.2) {\(K_{Y_2}\)};
\node[witness] (k2q) at (3.7,2.2) {\(K_{Y_q}\)};
\node[font=\scriptsize] at (-1.75,2.2) {\(\cdots\)};
\node[font=\scriptsize] at (3.05,2.2) {\(\cdots\)};
\foreach \target in {k11,k12,k1q}{\draw (t1)--(\target);}
\foreach \target in {k21,k22,k2q}{\draw (t2)--(\target);}
\draw[inclusion] (h)--(t1);
\draw[inclusion] (h)--(t2);
\node[panel,fit=(t1)(k11)(k12)(k1q)] (p1) {};
\node[panel,fit=(t2)(k21)(k22)(k2q)] (p2) {};
\node[font=\small,above=3pt of p1] {First private chart};
\node[font=\small,above=3pt of p2] {Second private chart};
\node[font=\small,below=7pt of h] {\(K_{X_i}\cap K_{Y_j}=H\quad\text{for every }i,j\)};
\end{tikzpicture}
\caption{The subgroup relations supplied by \Cref{c:failed-pattern}.
In the first chart, the witnesses on one parameter line satisfy
\(\Sigma_1=\langle X_i,X_j\rangle\) for every \(i\ne j\),
\(\operatorname{rk}\Sigma_1=\operatorname{rk}X_i+1\), and
\(T_1=K_{\Sigma_1}=K_{X_i}\cap K_{X_j}\); each \(K_{X_i}\) covers \(T_1\).
The second chart has the analogous relations with \(Y_j,\Sigma_2,T_2\).
Solid segments denote these covers, and dashed segments denote inclusions,
with intervening labels omitted. The diagram is schematic in rank and
shows only the stated relations.}
\label{c:pattern-figure}
\end{figure}


Assign each failed chart the ancestor/category/case of one homogeneous
pattern extracted by \Cref{c:failed-pattern}.
Suppose the exclusion argument for each such choice
prohibits either one such pattern or a pair in distinct charts. In the
first case no chart can fail in that case; in the second at most one can.
Thus the number of charts lost in a round is at most the number of cases
requiring two charts. These are the group-theoretic failure colors; the
pullback-flat colors used inside Hales--Jewett do not enlarge the
chart-loss count. Retaining one successful flat in every surviving chart
provides the anchors for the next round.

If $F$ is successful after the last round, every
$W\subsetneq X\subseteq F$ has $K_X\ge K_F$ and therefore avoids every
completed-category stabilizer relative to every ancestor tested.
The local conclusions consequently apply to all comparable candidates
in $F$, before a next path node is chosen. \Cref{c:ancestor-socles} makes
their simple groups actual and nested. This is the property used for
simultaneous root and tensor protection in \Cref{w:selection}.

\begin{proposition}[Finite schedule of simultaneous exclusions]\label{c:pruning-schedule}
Suppose that the local exclusions below prohibit their prescribed failed patterns with absolute case and chart-loss bounds. For a bounded path length, finitely many rounds then produce a flat of final increment \(r\) such that the subgroup label of every proper extension in it is contained in none of the prescribed proper full stabilizers relative to any high-type ancestor. The next local simple has the corresponding transitivity or full-spanning properties relative to all those ancestors.
\end{proposition}

\begin{proof}
Use a simultaneous initial round for all initial classical subspace exclusions relative to any classical ancestors (the ones allowing a more-than-height requirement); their pattern proofs need no earlier anchor. The remaining categories of the structure arguments can then be ordered as there, requiring only bounded line/orbit or independent-list bounds: e.g. direct line systems and then other direct systems, prime fields then ordered factors for linear type, partial ordered factors then prime fields then paired idempotents and remaining ordered factors for isometry types, elementary subgroups then tensor systems; in alternating ancestors subsets and block partitions then affine, Cartesian and diagonal structures. Subcases rendered impossible by an already forbidden structure need no round of their own. There are only boundedly many categories/ancestor instances altogether for bounded path length. We thus obtain a pruned flat of final increment \(r\) from \(W\); all labels at proper extensions within it avoid the exclusions, giving the needed next almost simple socles by the structural arguments, and giving transitivity or natural full spanning as applicable relative also to earlier ancestors. Selection of which extension to take (additional closure-field tensor/root protections within a same-characteristic run) uses \Cref{w:selection}. If the current parent instead is of bounded Lie rank, use \Cref{w:bounded-exclusion} (same graph pattern with a bounded independent list and no prior anchor, relative just to this parent).
\end{proof}

\begin{proposition}[Order of the finite parameter choices]\label{c:parameter-order}
The test-lattice parameters can be chosen in an order that meets all the pattern requirements and all the later path requirements. In particular the initial height-dependent pattern requirement does not create a circular dependence between \(m,n,N\). Once the local group results are applied, these choices force either the bounded-type contradiction or a path longer than the protected-path bound.
\end{proposition}

\begin{proof}
Here is one parameter order, always finite. Fix the high thresholds, the resulting bounded-type constants, and the absolute bound on protected high-type path lengths from the group arguments. Choose \(D\) exceeding this bound by several stages (allowing the displaced point start and a stage for the bounded-type contradiction). Choose \(q\), a power of 7 for instance, above all absolute line-count requirements. Choose \(r\) as in the final flag selection, large for bounded-type requirements and the exponential-gap union bound over at most \(D\) steps, so the path ranks are at most \(Dr\). Choose larger intermediate trial increments backwards from \(r\), the increment after the first round thus a constant independent of \(m,n\): each preceding increment merely must allow the Hales--Jewett bound for the next bounded-slot round, at any path rank up to \(Dr\). Choose \(m\) above the ranks of these trials, then \(n\) sufficiently larger, including for residual interval length below all of them and bounded required spans. Now \(N\) can allow the first-round Hales--Jewett bound including the larger initial slot requirement depending on \(n\), and any no-anchor bounded-type trial. Finally take \(g\) sufficiently large for the failed-chart bounds at every stage. Thus the separate arguments imply contradiction on this finite lattice: either a bounded-type parent gives the direct contradiction, or too long a protected path is forced.
\end{proof}
